% =============================================================================
%  Experiência profissional
% =============================================================================
\secao{Experiência Profissional}

\entrada{Consultor DBA (Engenharia de Sistemas)}{eSPap, Entidade de Serviços Partilhados da Administração Pública}{Regime híbrido}{04/2026 – atual}
\begin{itens}
  \item Administração de instâncias e bases de dados SQL Server (versões 2008 a 2025) num ambiente on-premises de grande escala da Administração Pública, integrado numa equipa de 3 DBA.
  \item Administração de bases de dados em Azure SQL Database e Azure SQL Managed Instance.
  \item Manutenção preventiva: índices e estatísticas, verificação de integridade (DBCC CHECKDB) e gestão de capacidade.
  \item Monitorização e análise de desempenho com DMVs, Extended Events e Profiler, incluindo o diagnóstico de queries lentas.
  \item Gestão de cópias de segurança e recuperação de dados, e participação na migração das instâncias de Always On para configuração standalone.
  \item Administração dos servidores dedicados de SSIS, SSAS e SSRS: instalação, manutenção, disponibilidade, gestão de acessos e conectividade.
  \item Gestão de segurança e acessos (logins, roles e permissões), automação de tarefas com SQL Server Agent e T-SQL, verificações pontuais com PowerShell, migrações e upgrades de versão, e documentação de procedimentos.
\end{itens}

\entrada{Assistente Técnico, Sistemas de Informação}{ULS Tâmega e Sousa, E.P.E.}{Penafiel}{06/2019 – 03/2026}
\begin{itens}
  \item Levantamento e estruturação dos 17 indicadores de desempenho contratualizados do projeto-piloto CRI-SM (Saúde Mental), com mapeamento de cada indicador às fontes de dados hospitalares e à plataforma institucional de BI.
  \item Desenvolvimento de um protótipo de BI, do data warehouse em SQL Server ao modelo dimensional e medidas DAX em Power BI, com painéis de acompanhamento do acesso a consultas, apresentado ao Conselho de Gestão do CRI-SM.
  \item Apoio administrativo em contexto hospitalar (2019 a 2024), com utilização diária dos sistemas de informação clínicos e administrativos.
\end{itens}

\entrada{Chefe de Departamento}{Grupo Ricon SA}{Famalicão}{2012 – 2019}
\begin{itens}
  \item Coordenação de uma equipa de 10 pessoas nos setores de acabamento, embalagem e expedição, incluindo o planeamento logístico dos fluxos de exportação.
  \item Controlo de inventário, stocks e economato do departamento.
  \item Suporte técnico de 1.ª linha aos utilizadores e às infraestruturas informáticas administrativas.
\end{itens}

\entrada{Assistente Comercial}{Sanvitêxteis, Indústria de Vestuário, Lda.}{Penafiel}{2002 – 2012}
\begin{itens}
  \item Gestão de clientes, fornecedores e transportadores, aquisição de matérias subsidiárias e logística de exportação.
\end{itens}
